\chapter{Multi-assinaturas}
\label{chapter:multisignatures}
% Implementação revista para esta edição contra a revisão Monero
% 3646f648db57f60cca86430e25a635d19fa9b92a.

Uma carteira Monero normal tem uma chave privada de vista e uma chave privada
de gasto. Quem conhece a chave de gasto inteira pode autorizar uma
transacção. Numa carteira de multi-assinatura, essa autoridade é repartida:
escolhem-se \(N\) participantes e um limiar \(M\), e pelo menos \(M\) deles
têm de cooperar para gastar. Um grupo 2-de-3, por exemplo, continua a operar
se um participante estiver indisponível, mas nenhum participante decide
sozinho.

Este resultado não é obtido por um tipo especial de saída. Na lista de
blocos, o endereço partilhado tem o formato de um endereço Monero e a
transacção final tem o formato de uma transacção RingCT vulgar. A política
\(M\)-de-\(N\) é aplicada pelas chaves, pelos dados conservados nas carteiras e
pelo protocolo interactivo que
produz a assinatura. A lista de blocos não anuncia se o dinheiro pertence a
uma pessoa, a uma empresa ou a um grupo de co-signatários.

As observações sobre a implementação neste capítulo referem-se à revisão do
código Monero datada de 14 de Agosto de 2026
\cite{monero-source-v16}. A descrição separa deliberadamente três coisas:
o que essa revisão implementa, o que o consenso verifica e o que continua
apenas conceptual. Quando dissermos que uma operação «existe», queremos dizer
que há código e uma interface de carteira para a executar nessa revisão, não
apenas uma construção plausível num artigo.

\section{Estado da implementação}
\label{sec:multisig-implementation-status}

\begin{center}
\small
\begin{tabular}{|>{\raggedright\arraybackslash}p{4.1cm}|
  >{\raggedright\arraybackslash}p{3.0cm}|
  >{\raggedright\arraybackslash}p{7.3cm}|}
\hline
\textbf{Componente} & \textbf{Estado} & \textbf{Significado} \\
\hline
Troca de chaves \(M\)-de-\(N\) & Implementada, experimental &
A carteira constrói uma chave de vista comum e uma chave de gasto repartida,
com um máximo de 16 participantes. \\
\hline
Recepção e sincronização & Implementadas, experimentais &
Os participantes detectam as mesmas saídas e trocam imagens de chave parciais
e valores aleatórios públicos. \\
\hline
Assinatura de despesas & Implementada, experimental &
A carteira constrói assinaturas CLSAG distribuídas e transacções
\texttt{RCTTypeBulletproofPlus}. \\
\hline
Verificação pelo consenso & Implementada como RingCT normal &
Os nós verificam a transacção final; não verificam uma política
\(M\)-de-\(N\) separada nem aprendem o número de signatários. \\
\hline
Sistema MMS & Implementado na carteira de linha de comandos &
Organiza e transporta as mensagens do protocolo; não altera a criptografia
nem o consenso. \\
\hline
Políticas em árvore & Não implementadas &
São uma extensão conceptual conservada para explicar propostas de capítulos
posteriores, não uma função da carteira Monero. \\
\hline
Multi-assinatura FCMP++ ou Carrot & Primitivas experimentais, sem integração &
O ramo FCMP++ de \texttt{monero-oxide} contém algoritmos SAL com FROST; a
carteira Monero examinada continua a assinar CLSAG e não oferece um fluxo
FCMP++ multi-assinatura. \\
\hline
\end{tabular}
\end{center}

A funcionalidade implementada continua marcada como \emph{experimental} e
está desactivada por omissão. A carteira de linha de comandos exige
\begin{verbatim}
set enable-multisig-experimental 1
\end{verbatim}
antes das operações correspondentes. O aviso não é decorativo: um erro de
protocolo, uma cópia de segurança inadequada ou um participante malicioso
pode bloquear fundos; certas falhas podem ainda expor segredos. A
multi-assinatura deve, por isso, ser tratada como um protocolo criptográfico
interactivo, e não como autenticação convencional com várias palavras-passe.

\section{Modelo e objectivos de segurança}

Representaremos os participantes por
\(\mathcal P=\{P_1,\ldots,P_N\}\). A política é válida quando
\begin{align*}
 1\leq M\leq N.
\end{align*}
Qualquer subconjunto com pelo menos \(M\) participantes deve conseguir
assinar; um subconjunto com menos de \(M\) não deve conseguir reconstruir a
chave de gasto nem produzir uma assinatura válida.

Há três propriedades diferentes que convém não confundir:
\begin{itemize}
 \item \emph{confidencialidade}: menos de \(M\) participantes não aprendem a
 chave de gasto completa;
 \item \emph{autorização}: menos de \(M\) participantes não conseguem criar
 a resposta secreta necessária à assinatura;
 \item \emph{disponibilidade}: a política tolera a ausência de até \(N-M\)
 participantes.
\end{itemize}
A última propriedade não obriga ninguém a colaborar. Um co-signatário pode
recusar uma proposta, apagar os seus dados ou abandonar a sessão do protocolo.
Multi-assinatura reduz o número de participantes indispensáveis; não garante
a cooperação dos participantes disponíveis.

Todos os membros aprendem a chave privada de vista comum. Conseguem, portanto,
detectar as saídas recebidas pelo endereço, descodificar os respectivos
montantes e observar as recepções da carteira. Depois de trocarem as imagens
de chave parciais, conseguem também acompanhar quais dessas saídas foram
gastas. O limiar protege a autoridade de \emph{gasto}, não a privacidade
interna entre os membros do grupo.

\section{Comunicação entre co-signatários}
\label{sec:communicating}

A criação da carteira e a assinatura são protocolos com várias mensagens.
Cada participante tem de saber quais mensagens pertencem à mesma sessão, quem
as enviou e em que ronda devem ser usadas. As mensagens de troca de chaves da
implementação têm assinaturas de autenticação, mas a primeira ronda contém um
segredo auxiliar que será comum ao grupo. É, por isso, necessário um canal
autenticado e confidencial para transportar os ficheiros ou cadeias de texto.

Autenticar apenas o nome do contacto não basta. Antes de receber fundos, os
participantes devem comparar por um canal independente:
\begin{itemize}
 \item a lista e a ordem canónica dos signatários;
 \item os valores \(M\) e \(N\);
 \item o endereço partilhado final;
 \item a chave pública de gasto e a chave pública de vista que lhe dão origem.
\end{itemize}
Esta confirmação posterior à troca de chaves detecta mensagens omitidas ou
substituídas e carteiras que acabaram, por engano, em contas diferentes. A
implementação inclui uma ronda final precisamente para assinar e comparar as
chaves públicas resultantes.

Os formatos actuais começam por \texttt{MultisigxV2R1} na primeira ronda e
\texttt{MultisigxV2Rn} nas rondas seguintes. O código consultado rejeita os
formatos V1, considerados inseguros. Forçar a aceitação de um conjunto
incompleto de mensagens pode produzir uma conta inconsistente ou tornar os
restantes participantes dependentes de uma chave conhecida por um só membro.

\section{Agregação das chaves de gasto}
\label{sec:key-aggregation}

Antes de construir uma política \(M\)-de-\(N\), vejamos um problema mais
pequeno. Alice escolhe \(a\in_R\mathbb Z_l\) e publica \(A=aG\); Bob escolhe
\(b\in_R\mathbb Z_l\) e publica \(B=bG\). A primeira ideia é somar as chaves:
\begin{align*}
 K^s=A+B=(a+b)G.
\end{align*}
Alice conhece uma parcela e Bob conhece a outra. Parece uma chave 2-de-2.

\subsection{A soma ingénua}
\label{sec:naive-key-aggregation}
\label{subsec:drawbacks-naive-aggregation-cancellation}

Suponhamos, porém, que Alice publica \(A\) primeiro. Bob pode escolher uma
chave pública \(T=tG\) que controla sozinho e anunciar
\begin{align*}
 B=T-A.
\end{align*}
A soma anunciada é então
\begin{align*}
 A+B=T,
\end{align*}
cuja chave privada \(t\) é conhecida por Bob. As assinaturas do traslado não
corrigem esta relação algébrica. Este é um ataque de cancelamento de chaves.

\subsection{Coeficientes de agregação}
\label{sec:robust-key-aggregation}

A defesa usada pela implementação associa um coeficiente a cada chave e ao
conjunto completo de chaves \cite{MRL-0009-multisig}. Sejam
\begin{align*}
 \mathbb B=(B_1,\ldots,B_t)
\end{align*}
as chaves componentes distintas, em ordem canónica. Para cada \(B_j\),
calcula-se
\begin{align*}
 \lambda_j
   =\mathcal H_n(B_j,B_1,\ldots,B_t,T_{\rm agg}),
 \qquad
 T_{\rm agg}=\texttt{Multisig\_key\_agg},
\end{align*}
e a chave pública de gasto é
\begin{align}
 K^s=\sum_{j=1}^{t}\lambda_jB_j.
 \label{eq:multisig-weighted-aggregate}
\end{align}
Quem conhece \(b_j\), onde \(B_j=b_jG\), conserva como parcela privada
\begin{align*}
 x_j=\lambda_jb_j.
\end{align*}
A dependência de \(\lambda_j\) em todo o conjunto impede que o último
participante escolha uma chave que cancele as anteriores sem resolver um
problema criptográfico. A ordenação é essencial: sem ela, duas carteiras com
as mesmas chaves poderiam calcular agregados diferentes.

\section{Construção de uma política M-de-N}
\label{sec:m-of-n}

A implementação actual obtém o limiar através de trocas Diffie--Hellman.
Definamos
\begin{align}
 q=N-M+1.
 \label{eq:multisig-share-group-size}
\end{align}
Cada chave privada componente é conhecida por um subconjunto de \(q\)
participantes. Há
\begin{align*}
 \binom Nq=\binom N{M-1}
\end{align*}
componentes distintas, e cada participante conserva
\begin{align*}
 \binom{N-1}{q-1}=\binom{N-1}{N-M}
\end{align*}
delas.

Por que funciona esta escolha? Se \(N-M\) participantes faltarem, restam
\(M\). O complemento de qualquer subconjunto com \(q=N-M+1\) elementos tem
apenas \(M-1\) elementos. Logo, os \(M\) presentes intersectam todos os
subconjuntos de \(q\) membros e, em conjunto, possuem todas as parcelas da
chave agregada. Com apenas \(M-1\) participantes pode faltar pelo menos uma
parcela.

\subsection{Chaves-base e chave de vista}

Cada participante \(P_i\) começa com dois escalares independentes:
\begin{align*}
 a_i&\in_R\mathbb Z_l, & u_i&\in_R\mathbb Z_l.
\end{align*}
A chave-base pública \(A_i=a_iG\) autentica as mensagens e entra nas trocas
Diffie--Hellman. O valor auxiliar \(u_i\) é enviado aos restantes membros na
primeira ronda, através do canal confidencial. Depois de ordenar as
contribuições, todos calculam a mesma chave privada de vista
\begin{align*}
 k^v=\mathcal H_n(u_1,\ldots,u_N),
 \qquad K^v=k^vG.
\end{align*}
A função hash inclui a serialização e a separação de domínio definidas pelo
formato da implementação. A notação abrevia esses pormenores para mostrar a
relação essencial.

Partilhar \(k^v\) é deliberado. Cada carteira tem de reconhecer as mesmas
saídas e calcular a parte comum das respectivas chaves privadas. Quem não
deve ver os montantes e destinatários internos não deve pertencer a esta
carteira.

\subsection{Segredos partilhados por subconjuntos}
\label{sec:smaller-thresholds}

Para um subconjunto \(S\subseteq\mathcal P\) com \(|S|=q\), as rondas
Diffie--Hellman produzem conceptualmente
\begin{align*}
 D_S=8^{q-1}\left(\prod_{i\in S}a_i\right)G,
 \qquad b_S=\mathcal H_n(D_S).
\end{align*}
O factor \(8\) em cada passo limpa o cofactor da curva. Cada membro de \(S\)
consegue chegar a \(D_S\) a partir da sua chave privada e dos pontos
recebidos; quem está fora de \(S\) não aprende \(b_S\).

Os pontos \(B_S=b_SG\) formam a lista \(\mathbb B\) da
equação \eqref{eq:multisig-weighted-aggregate}. Depois de calcular os
coeficientes de agregação, cada membro substitui a parcela que conhece por
\begin{align*}
 x_S=\lambda_Sb_S.
\end{align*}
A chave pública final é
\begin{align*}
 K^s=\sum_{S:\,|S|=q}\lambda_SB_S.
\end{align*}
Uma parcela conhecida por dois signatários continua a entrar uma única vez na
soma. Durante a assinatura, a carteira conserva o conjunto das parcelas já
incluídas para impedir duplicações.

\subsection{Rondas e confirmação do endereço}

O número de rondas criptográficas de troca de chaves é
\begin{align*}
 N-M+1=q.
\end{align*}
Segue-se uma ronda de confirmação, pelo que a configuração completa exige
\begin{align*}
 N-M+2
\end{align*}
rondas. Uma política \(N\)-de-\(N\) precisa de uma ronda de troca e uma de
confirmação; baixar o limiar acrescenta trocas Diffie--Hellman.

Em cada ronda posterior à primeira, um participante aplica a sua chave-base a
pontos cuja lista de origens ainda não o contém. A carteira acompanha essas
origens para que cada segredo final corresponda a exactamente \(q\) membros.
Na última troca reúne os pontos públicos finais, elimina duplicados, calcula
os coeficientes da agregação e actualiza as parcelas privadas locais. Na ronda
de confirmação, todos assinam as mesmas chaves públicas de gasto e de vista.

Só depois dessa confirmação o endereço deve receber fundos. Uma diferença de
um byte numa mensagem pode dar uma carteira válida mas diferente, o que é uma
forma particularmente pouco interessante de descobrir que a verificação foi
saltada.

\subsection{Exemplo 2-de-3}
\label{sec:n-1-of-n}

Consideremos Alice, Bob e Carlos, com \(M=2\) e \(N=3\). Temos \(q=2\), pelo
que são criados três segredos:
\begin{align*}
 b_{AB}&=\mathcal H_n(8a_Ab_BG),\\
 b_{AC}&=\mathcal H_n(8a_Aa_CG),\\
 b_{BC}&=\mathcal H_n(8a_Ba_CG).
\end{align*}
Depois da ponderação, Alice conhece \(x_{AB}\) e \(x_{AC}\), Bob conhece
\(x_{AB}\) e \(x_{BC}\), e Carlos conhece \(x_{AC}\) e \(x_{BC}\). A chave de
gasto é
\begin{align*}
 K^s=(x_{AB}+x_{AC}+x_{BC})G.
\end{align*}

Alice e Bob cobrem as três parcelas: Alice fornece \(x_{AC}\), Bob fornece
\(x_{BC}\), e um deles inclui \(x_{AB}\). O mesmo acontece com os pares
\((A,C)\) e \((B,C)\). Um participante isolado conhece apenas duas parcelas e
não consegue formar a chave completa.

Nos extremos, uma política \(N\)-de-\(N\) tem \(q=1\): cada participante
possui uma parcela exclusiva. Uma política 1-de-\(N\) tem \(q=N\): todos
derivam a mesma parcela e qualquer um consegue usá-la. Portanto, a política
1-de-\(N\) não oferece protecção contra um membro que decida gastar sozinho.

O número combinatório de parcelas cresce depressa. O código limita \(N\) a
16 participantes, um limite prático e não uma propriedade matemática do
protocolo. A própria implementação recomenda uma construção do género FROST
se esse limite algum dia tiver de crescer.

\section{Recepção de saídas e imagens de chave}
\label{sec:recalculating-key-images-multisig}

Como todos conhecem \(k^v\), todas as carteiras conseguem detectar uma saída
recebida. Para uma saída pública \(P\), escrevamos a respectiva chave privada
sob a forma
\begin{align}
 x=d+\sum_{j=1}^{t}x_j,
 \label{eq:multisig-output-secret}
\end{align}
onde \(d\) reúne a derivação comum obtida com a chave de vista e, quando
aplicável, o termo do sub-endereço; os \(x_j\) são as parcelas distintas da
chave de gasto. A imagem de chave é
\begin{align*}
 I=x\mathcal H_p(P)
   =d\mathcal H_p(P)+\sum_{j=1}^{t}x_j\mathcal H_p(P).
\end{align*}

Cada carteira calcula imagens parciais
\begin{align*}
 I_j=x_j\mathcal H_p(P)
\end{align*}
para as parcelas que possui. Depois de importar informação suficiente dos
outros signatários, reúne as imagens parciais distintas e acrescenta uma só
vez o termo comum \(d\mathcal H_p(P)\). Assim obtém a mesma imagem \(I\) que
uma carteira de uma só chave produziria.

Este passo é necessário mesmo antes de gastar. Uma carteira precisa da imagem
de chave completa para perguntar se a saída já apareceu como entrada na lista
de blocos. Ver a saída com a chave de vista informa que ela foi recebida; não
informa, por si só, se outro conjunto de \(M\) signatários já a gastou.

A operação de exportação de informação multi-assinatura inclui, para cada
saída, as imagens parciais e pares públicos de valores aleatórios que poderão
ser usados numa tentativa de assinatura. O ficheiro é cifrado com material
derivado da chave de vista comum e ligado à conta e ao signatário. Ao importar
dados de pelo menos \(M\) fontes contando a própria carteira, as imagens de
chave podem ser reconstruídas e o estado de gasto actualizado.

Uma resposta de um nó remoto não fiável acerca do estado de gasto pode ser
falsa. A criptografia valida a imagem de chave; não transforma o nó que
responde numa testemunha honesta da lista de blocos.

\section{Assinaturas de Schnorr com limiar}
\label{sec:threshold-schnorr}

O mecanismo de resposta é mais fácil de ver numa assinatura de Schnorr sem
anel. Suponhamos que a chave privada de uma saída está repartida como
\begin{align*}
 x=\sum_{i\in\mathcal A}x_i,
 \qquad X=xG,
\end{align*}
onde \(\mathcal A\) contém exactamente as parcelas distintas cobertas pelos
signatários escolhidos.

Cada participante escolhe um valor aleatório
\(\alpha_i\in_R\mathbb Z_l\) e publica \(L_i=\alpha_iG\). O ponto agregado e o
desafio são
\begin{align*}
 L&=\sum_{i\in\mathcal A}L_i,\\
 c&=\mathcal H_n(m,X,L).
\end{align*}
Cada participante devolve
\begin{align*}
 s_i=\alpha_i-cx_i,
\end{align*}
e a resposta final é
\begin{align*}
 s=\sum_{i\in\mathcal A}s_i.
\end{align*}
O verificador recompõe
\begin{align*}
 sG+cX
 &=\sum_i(\alpha_i-cx_i)G+c\sum_ix_iG\\
 &=\sum_i\alpha_iG=L
\end{align*}
e confirma o desafio. Nenhum participante precisou de reconstruir \(x\) num
único lugar.

O perigo clássico está na reutilização de \(\alpha_i\). Se o mesmo valor
aparecer em desafios distintos \(c\) e \(c'\), observam-se
\begin{align*}
 s_i&=\alpha_i-cx_i, &
 s'_i&=\alpha_i-c'x_i,
\end{align*}
e qualquer co-signatário calcula
\begin{align*}
 x_i=(s_i-s'_i)(c'-c)^{-1}.
\end{align*}
Um valor aleatório de assinatura é de uso único. Fazer uma cópia de segurança
desse valor e reutilizá-la não é prudência; é guardar duas metades de uma
equação que revela a chave.

Protocolos simples de duas rondas têm ainda ataques nos quais o último
participante escolhe o seu ponto aleatório depois de observar os restantes e
manipula muitos desafios relacionados \cite{cryptoeprint:2018:417}. Uma
solução histórica é comprometer primeiro os pontos e revelá-los depois. A
implementação actual segue uma solução do género MuSig2: usa dois valores
aleatórios independentes por participante e liga a combinação ao traslado
completo \cite{musig2}.

\subsection{Dois valores aleatórios por signatário}

Para cada entrada e tentativa, o signatário \(i\) prepara
\begin{align*}
 \alpha_{i,0},\alpha_{i,1}\in_R\mathbb Z_l
\end{align*}
e os pontos correspondentes necessários às duas bases da CLSAG. Seja
\(\tau\) o traslado que contém o anel, os compromissos, a compensação, a
mensagem, as imagens, as respostas simuladas, a posição real e as dimensões
do anel. Depois de somar separadamente os pontos públicos dos participantes,
\((L_0,L_1)\) e \((R_0,R_1)\), a carteira calcula
\begin{align*}
 b=\mathcal H_n(T_{\rm ms},\tau,L_0,L_1,R_0,R_1),
\end{align*}
com separação de domínio própria. O valor secreto combinado é
\begin{align}
 \alpha_i=\alpha_{i,0}+b\alpha_{i,1}.
 \label{eq:multisig-two-nonce}
\end{align}
Mudar a transacção, o anel, uma resposta simulada ou os pontos de qualquer
participante muda \(b\). O último signatário deixa de poder escolher
livremente o valor aleatório agregado depois de conhecer o traslado que o
determina.

\section{Assinatura CLSAG distribuída}
\label{sec:rcttypebulletproof2-multisig}

Uma transacção corrente usa CLSAG para provar que a entrada real está numa
posição desconhecida de um anel e que os compromissos de montante se
equilibram. A prova de intervalo Bulletproof+ trata os montantes das saídas;
ela não substitui a CLSAG que autoriza cada entrada.

Para a posição real, escrevamos \(x\) para a chave privada da saída e \(z\)
para o segredo do compromisso à diferença entre a entrada real e o
pseudo-compromisso. A CLSAG combina os dois testemunhos com coeficientes
públicos \(\mu_P\) e \(\mu_C\):
\begin{align*}
 w=\mu_Px+\mu_Cz.
\end{align*}
As imagens agregadas correspondentes são construídas a partir da imagem de
chave \(I=x\mathcal H_p(P)\) e da imagem de compromisso \(D\). Os restantes
membros do anel recebem respostas simuladas, como numa CLSAG normal.

Na versão distribuída, o primeiro signatário da tentativa inclui na sua
resposta:
\begin{itemize}
 \item as parcelas locais da chave de gasto ainda não usadas;
 \item o termo comum derivado da chave de vista e o ajuste de sub-endereço;
 \item o segredo do compromisso à diferença, ponderado por \(\mu_C\);
 \item a sua componente aleatória combinada da
 equação \eqref{eq:multisig-two-nonce}.
\end{itemize}
Cada signatário posterior acrescenta apenas as parcelas de gasto que ainda não
constam do conjunto \texttt{signing\_keys}, além da sua componente aleatória.
Quando todas as parcelas distintas estão cobertas, a resposta da posição real
é a mesma que seria obtida com o testemunho completo \(w\). O desafio inicial
da volta CLSAG e a resposta real são então inseridos numa assinatura CLSAG
ordinária.

O coeficiente \(b\) dos dois valores aleatórios inclui o anel, os
compromissos, a compensação, a mensagem, as imagens \(I\) e \(D\), todas as
componentes públicas aleatórias, as respostas falsas, a posição real e as
dimensões do anel. A finalidade desta lista extensa é simples: não deixar um
pormenor relevante da tentativa fora do compromisso criptográfico.

Na revisão estudada, o construtor multi-assinatura aceita as versões RingCT
activas que usam CLSAG. A saída produzida tem o tipo
\texttt{RCTType}\allowbreak\texttt{BulletproofPlus}. O nome antigo
\texttt{RCTType}\allowbreak\texttt{Bulletproof2}, conservado no rótulo desta
secção para não quebrar referências de outras edições, já não descreve o
resultado.

\section{Construção e assinatura da transacção}

Uma carteira individual pode escolher entradas, saídas, troco e taxa e assinar
de imediato. Uma carteira \(M\)-de-\(N\) tem de fazer os mesmos cálculos em
várias máquinas e obter respostas compatíveis. O processo separa
sincronização, proposta, assinatura e submissão.

\subsection{Sincronização da carteira}

Antes de propor uma transacção, os participantes trocam informação
multi-assinatura. A exportação contém as imagens de chave parciais e, por
saída, os pares públicos de valores aleatórios preparados para futuras
tentativas. A importação combina imagens, actualiza o estado das saídas e
regista quais signatários têm material disponível.

Exportar novamente substitui os valores aleatórios privados antigos. As
tentativas que dependiam deles deixam de poder ser concluídas. Este
comportamento ajuda a impor o uso único, mas significa que os ficheiros de uma
mesma sessão não devem ser misturados ao acaso.

\subsection{Proposta}

O proponente escolhe as entradas e destinos, calcula os pseudo-compromissos e
a prova Bulletproof+, e fixa a estrutura que todos irão verificar. Cria também
tentativas para as combinações de signatários que podem atingir o limiar com
o material importado.

Numa carteira 2-de-4, se Alice propõe e tem dados de Bob, Carlos e Diana, pode
preparar tentativas \(AB\), \(AC\) e \(AD\). Cada tentativa usa pontos
aleatórios próprios. Assim, a indisponibilidade de Bob não obriga a reconstruir
todo o pagamento, desde que outra tentativa ainda tenha valores privados
válidos.

Dois proponentes não devem construir propostas concorrentes com os mesmos
pares públicos. Uma resposta enviada para uma delas consome os respectivos
segredos, mesmo que a transacção nunca seja submetida. Reutilizar o par na
outra proposta abre a equação de extracção vista na secção anterior.

A chave privada da transacção é derivada de entropia nova e das imagens de
chave das entradas. Isto evita que propostas repetidas criem as mesmas chaves
de saída de uso único e tornem inacessíveis fundos destinados a transacções
diferentes.

\subsection{Respostas e submissão}

Ao receber o conjunto de transacções, cada carteira reconstrói a proposta e
verifica a estrutura, as entradas, os destinos, a taxa, o equilíbrio dos
compromissos e a prova de intervalo. O utilizador deve ainda confirmar os
destinos e montantes no seu próprio dispositivo. Uma estrutura
criptograficamente válida pode conter um endereço de destino incorrecto.

Se aceitar, o signatário acrescenta as suas respostas parciais a todas as
tentativas em que participa. Antes de expor essas respostas, a carteira apaga
o material aleatório guardado para as tentativas e grava o estado actualizado.
Esta ordem procura garantir que uma falha posterior ou uma segunda cópia do
ficheiro não provoque reutilização.

Quando uma tentativa reúne todas as parcelas necessárias, a assinatura CLSAG
fica completa. O último participante pode então importar o conjunto, verificar
a transacção final e submetê-la à rede. Os nós e os mineiros verificam uma
transacção RingCT normal; não conhecem os nomes, o limiar nem o número total de
co-signatários.

\subsection{Comunicação simplificada}
\label{sec:simplified-communication}

Em termos operacionais, uma sessão correcta segue esta sequência:
\begin{enumerate}
 \item todos sincronizam a carteira com a lista de blocos;
 \item cada signatário exporta informação nova uma vez;
 \item o proponente importa informação de um conjunto suficiente e cria a
 proposta;
 \item os participantes verificam e assinam a mesma proposta, sem gerar uma
 nova exportação a meio;
 \item um participante importa as respostas completas e submete a
 transacção;
 \item todos voltam a sincronizar e trocam informação nova antes da próxima
 despesa.
\end{enumerate}
O Sistema de Mensagens de Multi-assinatura, MMS, automatiza o transporte e o
estado desta troca de protocolo na carteira de linha de comandos
\cite{mms-manual,mms-project-proposal}. Automatizar reduz erros de operação,
mas não remove a necessidade de autenticar os contactos e confirmar a
transacção.

\section{Sequência mínima de comandos}
\label{sec:multisig-cli}

Os comandos seguintes não constituem um guia de instalação; fixam apenas a
correspondência entre o protocolo matemático e a interface que realmente o
executa. Cada participante começa com uma carteira nova e nunca usada. Em
cada carteira:
\begin{verbatim}
set enable-multisig-experimental 1
prepare_multisig
make_multisig <M> <R1-P1> [<R1-P2> ...]
exchange_multisig_keys <Rk-P1> [<Rk-P2> ...]
\end{verbatim}
O resultado de cada comando de troca é enviado aos restantes participantes.
Repete-se o último comando enquanto a carteira disser que falta outra ronda.
Matematicamente, esta sequência constrói \(k^v\), as parcelas \(x_S\) e
\[
 K^s=\sum_{S:\,|S|=N-M+1}x_SG.
\]
Todos confirmam que a carteira terminou com o mesmo endereço antes de esse
endereço receber fundos.

Para gastar uma saída, os signatários escolhidos actualizam primeiro as
imagens de chave e os pares públicos de valores aleatórios:
\begin{verbatim}
export_multisig_info info-Pi
import_multisig_info info-P1 [info-P2 ...]
\end{verbatim}
Os ficheiros exportados são trocados entre eles. O proponente fixa a mensagem
da transacção, os anéis, os compromissos, os destinos e a taxa:
\begin{verbatim}
transfer <endereco-destino> <montante>
\end{verbatim}
Numa carteira multi-assinatura, \texttt{transfer} escreve a proposta no
ficheiro \texttt{multisig\_monero\_tx}; não a difunde. Os co-signatários
passam a versão actualizada desse mesmo ficheiro entre si e executam
\begin{verbatim}
sign_multisig multisig_monero_tx
\end{verbatim}
até uma tentativa conter as \(M\) contribuições necessárias. Cada chamada
acrescenta as parcelas ainda ausentes à resposta CLSAG
\[
 s=\sum_i\alpha_i-c\!
   \left(\mu_P\left(d+\sum_i x_i\right)+\mu_Cz\right).
\]
Depois da última verificação, uma carteira difunde a transacção completa:
\begin{verbatim}
submit_multisig multisig_monero_tx
\end{verbatim}
Assim, a sequência operacional corresponde às equações anteriores:
trocar chaves para obter \(K^s\), trocar dados para obter \(I\) e os pontos
aleatórios, fixar o traslado, somar respostas e submeter a CLSAG resultante.
O MMS automatiza o transporte e o estado destas mensagens
\cite{mms-manual,mms-project-proposal}; não altera a matemática nem o
consenso.

Na revisão consultada aplicam-se ainda os seguintes limites:
\begin{itemize}
 \item a funcionalidade tem de ser activada explicitamente por ser
 experimental;
 \item o máximo é 16 signatários;
 \item as operações multi-assinatura indicadas não são suportadas por
 carteiras de hardware;
 \item mensagens de rondas antigas ou incompletas não devem ser forçadas;
 \item uma cópia de uma carteira já convertida não deve actuar como um novo
 participante independente.
\end{itemize}

Uma carteira clonada possui as mesmas parcelas e pode repetir material de
assinatura que a original julgava consumido. As cópias de segurança devem ser
entendidas como recuperação do mesmo participante, não como um participante
adicional no limiar.

\section{Políticas compostas não implementadas}
\label{sec:general-key-families}

A construção implementada é plana: \(M\) de um conjunto fixo de \(N\)
participantes. Não implementa políticas como «Alice e um de três auditores»
ou «dois departamentos, cada um representado por uma de duas pessoas».
Essas expressões formam árvores de limiares e pertencem, nesta edição, à
análise conceptual.

\subsection{Agregação conceptual por conjuntos de chaves}
\label{subsec:nesting-multisig-keys}

Conceptualmente, cada folha de uma árvore teria uma chave pública. Um nó
interno agregaria as chaves dos filhos segundo a sua própria regra de limiar,
e a raiz forneceria a chave pública de gasto. Uma extensão possível permitiria que
um «participante» de uma carteira exterior fosse ele próprio uma carteira
interior.

Esta ideia é conservada porque capítulos posteriores a usam ao discutir
serviços de garantia e grupos de moderadores. Não é, contudo, uma capacidade
da implementação revista. A troca de chaves Monero aceita um único par
\(M,N\) e uma lista plana de signatários; a interface não serializa uma
árvore arbitrária, não coordena os seus valores aleatórios e não prova que os
ramos interiores executaram as políticas anunciadas.

Certas políticas podem ser simuladas distribuindo dispositivos ou escolhendo
cuidadosamente as parcelas, mas isso altera os pressupostos de confiança e
pode aumentar a probabilidade de reutilização de valores aleatórios. Uma
implementação real de políticas compostas precisaria de um protocolo próprio,
análise de segurança, formatos de mensagens e testes de recuperação. Nada
disso está presente no código consultado.

\clearpage
\section{Nota sobre multi-assinatura com FCMP++}
\label{sec:multisig-fcmp-note}

\textbf{Resposta curta: ainda não está disponível na carteira Monero, mas já
existem primitivas experimentais no ramo de referência consultado.} FCMP++
separa a prova de pertença à cadeia da prova SAL de autorização e ligação
\cite{fcmp-plus-plus-paper}. Um coordenador pode construir o caminho de
pertença sem conhecer a chave de gasto inteira. A cooperação \(M\)-de-\(N\)
concentra-se na imagem de chave e na SAL.

Para uma saída de uma carteira de hierarquia antiga, escrevamos
\begin{align*}
 x&=d+\sum_{j\in\mathcal A}x_j,&
 I&=\mathcal H_{\rm KI}(K_o),\\
 L=xI&=dI+\sum_{j\in\mathcal A}x_jI.
\end{align*}
\(\mathcal A\) contém uma vez cada parcela coberta pelos signatários. A imagem
\(L\) é, portanto, a soma de imagens parciais; \(\mathcal H_{\rm KI}\) denota
a derivação de base exigida pela versão da saída. Não se inverte a chave
repartida. Esta linearidade torna a multi-assinatura viável.

A SAL prova também conhecimento de \(x\), da outra abertura
\(\widetilde y\) em \(\widetilde O=xG+\widetilde yT\), e do produto
\(z=xr\). Fixado o mesmo re-randomizador \(r\), podem formar
\begin{align*}
 z&=rd+\sum_{j\in\mathcal A}rx_j,&
 s_x&=\sum_j\alpha_j+ex,&
 s_z&=\sum_j\rho_j+ez.
\end{align*}
Nenhum signatário pode enviar \(rx_j\) sem máscara, pois \(r\) revelaria
\(x_j\). Envia apenas pontos públicos e respostas ocultas pelos valores
aleatórios \(\alpha_j,\rho_j\), usados uma vez. Um protocolo completo tem
ainda de construir conjuntamente \(P,A,B,R_O,R_P,R_L\), tratar a abertura em
\(T\) e resistir a valores duplicados e abortos maliciosos. As somas acima são
um esboço algébrico, não o protocolo distribuído da implementação de referência
nem uma análise de segurança.

\subsection{As duas primitivas no ramo de referência}

Na revisão \texttt{31c26d96eaadbba910ffe3613ad8b4cf9c598a93} do ramo
\texttt{fcmp++} de \texttt{monero-oxide}, o módulo SAL expõe, sob a
funcionalidade opcional \texttt{multisig}, dois algoritmos baseados em FROST
\cite{monero-oxide-fcmp-source}:
\begin{itemize}
 \item \texttt{SalLegacyAlgorithm} recebe uma chave \(x\) repartida, como nas
 carteiras históricas, e uma abertura \(y\) conhecida. Cada participante
 acrescenta parcelas interpoladas da imagem de chave e de \(xU\); a resposta
 FROST comum produz simultaneamente as relações em \(G\), \(U\) e
 \(\widetilde I\).

 \item \texttt{SalAlgorithm} destina-se à decomposição nova
 \(O=xG+yT\). O escalar \(x\), necessário para gerar a imagem de chave, é
 conhecido, enquanto \(y\), a componente em \(T\) que autoriza o gasto, é
 repartido com FROST sobre a base \(T\). A resposta partilhada torna-se
 \(s_y\); as restantes respostas SAL são construídas a partir dos dados comuns
 e de valores aleatórios ligados ao traslado.
\end{itemize}
Os dois algoritmos têm testes unitários que agregam parcelas e verificam a SAL
resultante. Isto é implementação de primitivas criptográficas, não integração
na carteira Monero nem activação pelo consenso: o módulo está condicionado por
uma opção de compilação e não define por si só criação de contas, transporte
de mensagens, recuperação, selecção de entradas ou comandos de utilizador.

Carrot acrescenta outra mudança. A hierarquia nova usa
\begin{align*}
 K^s=k_{\rm gi}G+k_{\rm ps}T
\end{align*}
e especifica que o segredo mestre não existe numa carteira
multi-assinatura \cite{carrot-specification}. A primitiva nativa acima conserva
\(k_{\rm gi}\), necessário à imagem de chave, como componente conhecida e
reparte \(k_{\rm ps}\), a componente de autorização em \(T\). Assim, uma vista
completa pode actualizar o estado gasto e o saldo sem reconstruir a parcela
privada que autoriza o gasto.

Na cadeia ver-se-iam apenas uma prova de pertença, uma SAL e \(L\), nunca
\(M\), \(N\) ou os signatários. No código Monero consultado, porém, o
construtor multi-assinatura termina numa CLSAG: não há comandos nem mensagens
que liguem a carteira às primitivas SAL distribuídas do ramo de referência.
Logo, \emph{as primitivas existem, mas a funcionalidade ainda não está
disponível na carteira}.

\clearpage
\section{Resumo}

Uma carteira Monero \(M\)-de-\(N\) transforma a chave de gasto numa soma
ponderada de parcelas Diffie--Hellman. Cada parcela é conhecida por
\(N-M+1\) participantes; por construção, qualquer grupo de \(M\) cobre todas
as parcelas e um grupo menor não as cobre. Os coeficientes ligados ao conjunto
completo de chaves impedem cancelamentos durante a agregação.

A chave de vista é comum a todos os membros. Para reconhecer despesas, as
carteiras trocam imagens de chave parciais; para autorizar uma entrada, somam
respostas parciais numa CLSAG. Dois valores aleatórios por signatário, ligados
ao traslado completo, defendem a assinatura de ataques adaptativos, e cada par
tem de ser usado uma única vez.

O resultado publicado é uma transacção RingCT normal. A complexidade fica nas
carteiras: rondas autenticadas e confidenciais, confirmação do endereço,
sincronização, verificação independente da proposta, destruição segura dos
valores aleatórios e cópias de segurança coerentes. É precisamente por essa
complexidade que a implementação consultada continua assinalada como
experimental.
